\documentclass[a4paper,11pt]{article}
\usepackage[utf8]{inputenc}
\usepackage[T1]{fontenc}
\usepackage[polish]{babel}
\usepackage{url}
\usepackage{geometry}
\usepackage{amsmath}

\geometry{
 a4paper,
 total={170mm,257mm},
 left=20mm,
 top=20mm,
}

\title{\textbf{Wniosek o Grant Obliczeniowy PLGrid}}
\author{Paweł Pozorski, Jakub Muszyński}
\date{\today}

\begin{document}

\maketitle
\begin{abstract}
\begin{center}
\begin{tabular}{@{}p{0.28\textwidth} p{0.65\textwidth}@{}}
\textbf{Tytuł projektu:} & Analiza wyjaśnialności multimodalnych i wielojęzycznych modeli językowych (MLLM) z wykorzystaniem wartości Shapleya \\[4pt]
\textbf{Słowa kluczowe:} & XAI, MLLM, Shapley Values, NLP, Audio processing, Large Language Models, Explainability
\end{tabular}
\vspace{10pt}

Celem projektu jest przeprowadzenie eksperymentalnej analizy wyjaśnialności (XAI) multimodalnych dużych modeli językowych (MLLM), przetwarzających zarówno tekst, jak i dźwięk. Badania stanowią kluczowy element pracy inżynierskiej realizowanej na Wydziale Matematyki i Nauk Informacyjnych Politechniki Warszawskiej.
\end{center}
\end{abstract}

\section*{Założenia projektu}

Projekt zakłada wykorzystanie autorskiego pakietu \texttt{mllm-shap} \footnote{\url{https://pawlo77.github.io/MLLM-Shap/}} do estymacji wartości Shapleya (SV) dla modeli typu \textit{audio-to-audio} oraz \textit{text-audio}. Głównym celem badawczym jest zrozumienie, w jaki sposób informacje pochodzące z różnych modalności (tekst vs dźwięk) oraz języków (angielski, hiszpański, francuski) wpływają na decyzje modelu w konwersacjach wieloturowych. Badania obejmą analizę wpływu modalności wejścia/wyjścia na ,,skupienie'' (\textit{focus}) modelu oraz ocenę stabilności metod aproksymacji SV (Monte Carlo, alokacja Neymana) w kontekście danych multimodalnych.

\section*{Metodyka badań}

Badania opierają się na metodach wyjaśnialności typu \textit{black-box}, traktując model jako wyrocznię (bez dostępu do gradientów). Do obliczania atrybucji wykorzystane zostaną wartości Shapleya. Ze względu na wykładniczą złożoność obliczeniową dokładnego SV ($O(2^n)$), w projekcie zastosowane zostaną zaawansowane metody aproksymacji oparte na próbkowaniu:

\begin{enumerate} \item \textbf{Estymatory Monte Carlo:} Warianty Standard oraz Limited (z wymuszeniem oceny koalicji pierwszego rzędu -- \textit{first-order omission}). \item \textbf{Metoda Complementary Contributions (CC) z alokacją Neymana:} Wykorzystanie \textit{stratified sampling} w celu redukcji wariancji estymatora przy ograniczonym budżecie obliczeniowym. \end{enumerate}

\noindent Eksperymenty zostaną przeprowadzone na przygotowanych podzbiorach danych \textbf{VoiceBench} oraz \textbf{InfinityInstruct} (warianty \textit{single-sentence}, \textit{multi-sentence}, \textit{multi-turn}), każdy z nich w okolicach $100$ rekordów. Budżet aproksymacji ustawiony zostanie na $4 n^2$ gdzie $n$ - ilość tokenów wejścia, dając maksymalnie $8,000$ requestów do modelu dla każdego wejścia. Badane modele to m.in. Qwen2.5-Omni-3B, Baichuan-Audio (7B) oraz LFM2-Audio-1.5B.

\section*{Oprogramowanie}

\begin{itemize} \item Język: Python 3.10+ \item Biblioteki ML/DL: PyTorch, Transformers (Hugging Face) \item Biblioteki pomocnicze: NumPy, Scikit-learn, Pandas \item Autorskie oprogramowanie: Pakiet \texttt{mllm-shap} (dostępne przez \texttt{pip install}). \end{itemize}

\section*{Uzasadnienie zapotrzebowania na zasoby}

Realizacja projektu wymaga znaczącej mocy obliczeniowej GPU ze względu na specyfikę metod wyjaśnialności opartych na wartościach Shapleya. Wyznaczenie atrybucji dla pojedynczej próbki danych wymaga wielokrotnego (rzędu setek lub tysięcy razy, w zależności od budżetu próbkowania) uruchomienia inferencji modelu dla różnych permutacji maskowania wejścia.

Zgodnie ze wstępnymi estymacjami przeprowadzonymi na lokalnej stacji roboczej (Nvidia RTX 4080), przetworzenie jednego zapytania do modelu zajmuje okolice $2$ s. Planowane eksperymenty są zatem wymagające, ponieważ zakładają:

\begin{itemize} \item Analizę wartości Shapleya w zależności od modalności i języka. Powoduje to konieczność wielokrotnej iteracji zbiorów danych. \item Użycie modeli o większej liczbie parametrów (np. Baichuan-Audio 7B), co wymaga kart GPU o dużej pamięci VRAM (klasy A100) dla efektywnego przetwarzania wsadowego (\textit{batching}) i uniknięcia błędów OOM (Out Of Memory). \end{itemize}

Wstępne szacunki zasobów obliczeniowych dla budżetu $4 n^2$ oraz planowanych eksperymentów są następujące:
\begin{itemize}
    \item \textbf{single\_\_sentence}: $100$ próbek x $2$ modalności wyjścia x $3$ modalności wejścia per model. Dokładna liczba zapytań do modelu to $102,768$.
    \item \textbf{multi\_\_sentence}: $100$ próbek x $4$ konfuguracje modalności wejścia. Dokładna liczba zapytań do modelu to $814,416$.
    \item \textbf{multi\_\_lingual}: $108$ próbek x $3$ kofiguracje modalności wejścia. Dokładna liczba zapytań do modelu to $317,664$.
\end{itemize}

Łącznie daje to około $1,234,848$ zapytań do modelu. Przy średnim czasie przetwarzania $2$ s na zapytanie, całkowity czas obliczeń wynosi około $2,469,696$ s, czyli około $29$ dni ciągłej pracy na pojedynczym GPU. Każdy eksperyment musi zostać wykonany na $3$ różnych modelach, co daje łącznie około $87$ dni pracy na pojedynczym GPU. Jest to wersja optymisstyczna, ponieważ liczba tokenów audio może być większa niż tekstowych, oraz prędkość inferencji może się różnić w zależności od modelu i modalności wejścia (wybrane $2$ są dla modelu LFM2-Audio-1.5B z wejściem tekstowym).

Dostęp do klastra obliczeniowego pozwoli na zrównoleglenie obliczeń dla różnych modeli i konfiguracji językowych, co jest niezbędne do ukończenia eksperymentów w ramach harmonogramu pracy dyplomowej.

\end{document}
